\Needspace{11\baselineskip}
\SourceLeaf{281}
\begin{center}
\textit{Comparing the numbers of classes in the individual genera of different orders.}
\end{center}

\Article{253}
The preceding theorem assumes that the order is the same and must not be extended to different orders. Thus, \textit{for example,}\ for the determinant $-171$ there are $20$ positive classes, which fall into four orders: the properly primitive order contains two genera, each comprising six classes; in the improperly primitive order two genera have four classes, two each; in the order derived from the properly primitive order of determinant $-19$ there is a single genus comprising three classes; finally, the order derived from the improperly primitive order of determinant $-19$ has a single genus consisting of one class. The negative classes behave in the same way. It is therefore worth investigating the general principle on which the connection between the numbers of classes in different orders depends. Suppose $K,L$ are two classes of the same (positive) order $O$ of determinant $D$, and $M$ a properly primitive class of the same determinant whose composition with $K$ gives $L$, such as can always be found by Article~251. In some cases $M$ may be the \textit{only} properly primitive class whose composition with $K$ produces $L$; in others several distinct properly primitive classes may have this property. Suppose in general that there are $r$ such properly primitive classes $M,M',M''\ldots M^{r-1}$, each of which, composed with $K$, gives the same class $L$, and denote their collection by $W$. Further, let $L'$ be another class of the order $O$ (different from $L$), and $N'$ a properly primitive class of determinant $D$ whose composition with $L$ gives $L'$, and denote the collection of classes $N'+M,N'+M',N'+M''\ldots N'+M^{r-1}$ (all properly primitive and distinct) by $W'$. It is then easy to see that $K$, composed with any class in $W'$, produces $L'$, so $W$ and $W'$ have no class in common. We also verify without difficulty that no properly primitive class outside the collection $W'$, when composed with $K$, produces $L'$. Likewise, if $L''$ is another class of the order $O$, different from the classes $L,L'$, there are $r$ properly primitive forms, distinct both from one another and from the forms $W,W'$, each of which, composed with $K$, produces $L''$. The same holds for all the remaining classes of the order $O$. Since every properly primitive (positive) class of determinant $D$, composed with $K$, produces a class of the order $O$, we easily conclude that if the total number of classes of the order $O$ is $n$, the total number of properly primitive (positive) classes of the same determinant is $rn$.

\SourceLeaf{282}We therefore have the general rule: if \(K\), \(L\) denote any classes of the order \(O\), and \(r\) the number of distinct properly primitive classes of the same determinant whose compositions with \(K\) each produce \(L\), the total number of classes in the properly primitive (positive) order is \(r\) times as large as the number of classes of the order \(O\).

Since the classes \(K\), \(L\) in the order \(O\) can be chosen entirely at will, we may also take identical classes. It is useful to take the class. which contains the simplest form of this order. Taking this for \(K\) and \(L\), then, reduces the matter to finding all the properly primitive classes whose compositions with \(K\) reproduce \(K\) itself. The following prepares the way for this.

\Needspace{8\baselineskip}
\Article{254}
\textsc{Theorem.} \textit{If \(F=(A,B,C)\) is the simplest form of the order \(O\) of determinant \(D\), and \(f=(a,b,c)\) a properly primitive form of the same determinant, this form \(f\) can represent the number \(AA\) if \(F\) is obtained by composing the forms \(f\), \(F\); and conversely, \(F\) is composed from itself and \(f\) if \(AA\) can be represented by \(f\).}

\textit{Proof.} I. If \(F\) is transformed into the product \(fF\) by the substitution \(p,\allowbreak p',\allowbreak p'',\allowbreak p''';\allowbreak \ q,\allowbreak q',\allowbreak q'',\allowbreak q'''\), Article 235 gives
\[
\begin{gathered}A(aq''q''-2bqq''+cqq)=A^3,\quad\text{so}\\ AA=aq''q''-2bqq''+cqq.\quad\textit{First part proved.}\end{gathered}
\]

II. \textit{Assume} that \(AA\) can be represented by \(f\), and denote the values of the indeterminates that give this representation by \(q''\). \(-q\), so that \(AA=aq''q''-2bqq''+cqq\), and set
\[
\begin{aligned}
q''a-q(b+B)&=Ap,\\ -qC&=Ap',\\ q''(b-B)-qc&=Ap''\\ -q''C&=Ap''',\\ q''a-q(b-B)&=Aq',\\ q''(b+B)-qc&=Aq'''
\end{aligned}
\]
Having done this, we easily confirm that \(F\) is transformed into the product \(fF\) by the substitution \(p,\allowbreak p',\allowbreak p'',\allowbreak p''';\allowbreak \ q,\allowbreak q',\allowbreak q'',\allowbreak q'''\). and is therefore composed from \(f\) and \(F\), provided all the numbers \(p\), \(p'\) etc. are integers. Now, by the description of the simplest form, \(B\) is either \(0\) or \(\frac12 A\), so \(\frac{2B}{A}\) is an integer; the same description shows that \(\frac{C}{A}\) is always an integer. Thus \(q'-p\), \(p'\), \(q'''-p''\), \(p'''\) are integers, and it only remains to prove that \(p\) and \(p''\) are integers. But we have
\SourceLeaf{283}
\[
pp+\frac{2pqB}{A}=a-\frac{qqC}{A},\quad p''p''+\frac{2p''q''B}{A}=c-\frac{q''q''C}{A}
\]
so if \(B=0\), we have
\[
pp=a-\frac{qqC}{A},\quad p''p''=c-\frac{q''q''C}{A}
\]
and therefore \(p,p''\) are integers; but if \(B=\frac12 A\), we have
\[
pp+pq=a-\frac{qqC}{A},\quad p''p''+p''q''=c-\frac{q''q''C}{A}
\]
from which it follows just as easily that \(p\) and \(p''\) are integers in this case too. We conclude that \(F\) is composed from \(f\) and \(F\). \textit{Second part proved.}

\Needspace{8\baselineskip}
\Article{255}
The problem is therefore reduced to finding all properly primitive classes of determinant \(D\) whose forms can represent \(AA\). Clearly \(AA\) can be represented by any form whose leading coefficient is either \(AA\) or the square of a divisor of \(A\). Conversely, if \(AA\) can be represented by a form \(f\), assigning to its indeterminates the values \(\alpha e\), \(\gamma e\) with greatest common divisor \(e\), the form \(f\) is transformed by the substitution \(\alpha,\beta,\gamma,\delta\) into a form whose leading coefficient is \(\frac{AA}{ee}\). This form is properly equivalent to \(f\) if \(\beta\), \(\delta\) are chosen so that \(\alpha\delta-\beta\gamma=1\). Thus any class whose forms can represent \(AA\) contains forms whose leading coefficient is \(AA\) or the square of a divisor of \(A\). The task is therefore to find all properly primitive classes of determinant \(D\) in which such forms occur. This is done as follows. Let \(a,a',a''\) etc. be all the (positive) divisors of \(A\); find all the values of the expression \(\surd D\ (\mathrm{mod.}\,aa)\) lying between \(0\) and \(aa-1\) inclusive, denoting them by \(b,b',b''\) etc., and set
\[
bb-D=aac,\quad b'b'-D=aac',\quad b''b''-D=aac''\text{ etc.}
\]
Denote the collection of forms \((aa,b,c)\), \((aa,b',c')\) etc. by \(V\). It is easy to see that any class of determinant \(D\) containing a form with leading coefficient \(aa\) must also contain a form from \(V\). Likewise find all forms of determinant \(D\) whose leading coefficient is \(a'a'\) and whose middle coefficient lies between \(0\) and \(a'a'-1\) inclusive, and denote their collection by \(V'\); similarly let \(V''\)
\SourceLeaf{284}
be the collection of such forms with leading coefficient \(a''a''\), and so on. Remove from \(V,V',V''\) etc. all forms that are not properly primitive, sort the remaining forms into classes, and if several belong to the same class, retain only one in each class. This gives all the required classes, and their number is to unity as the number of all properly primitive (positive) classes is to the number of classes in the order \(O\).

\textit{Example.} Let \(D=-531\), and let \(O\) be the positive order derived from the improperly primitive order of determinant \(-59\), whose simplest form is \((6,3,90)\), so \(A=6\). Here \(a,a',a'',a'''\) are \(1,2,3,6\); \(V\) contains the form \((1,0,531)\); \(V'\) the forms \((4,1,133)\), \((4,3,135)\); \(V''\) the forms \((9,0,59)\), \((9,3,60)\), \((9,6,63)\); finally \(V'''\) the forms \((36,3,15)\), \((36,9,17)\), \((36,15,21)\), \((36,21,27)\), \((36,27,35)\), \((36,33,45)\). But six of these twelve forms must be removed: the second and third from \(V''\), and the first, third, fourth and sixth from \(V'''\), all of which are derived forms. The six remaining forms are found to belong to distinct classes. In fact, the number of properly primitive (positive) classes of determinant \(-531\) is \(18\), and the number of improperly primitive (positive) classes of determinant \(-59\) (or the number of classes of determinant \(-531\) derived from them) is \(3\), so the former is to the latter as \(6\) to \(1\).
